\section{Good Fit, Model Restrictions, and Causal Identification}
\label{sec:identification-limits}
The preceding problems concern the connection between the model and the
recorded correlations. A further problem would remain even if that connection
were correct: matching the moments does not uniquely establish their cause.
Correct specification and identification can justify consistency under
appropriate estimation conditions; they do not guarantee finite-sample
unbiasedness. Conversely, misspecified equations can fit closely while
assigning the wrong structural meaning to their parameters.

\subsection{The Lesson of the 1970s IQ Debate}
\label{sec:1970s-precedent}
The disagreement is not new. Models fitted to overlapping IQ kinship data
in the 1970s yielded sharply different heritability estimates despite
apparently satisfactory fit. \citet{Loehlin1978} traced important differences
to assumptions about dominance, assortment, and twin environments. His
constructive recommendation was to identify the consequential assumptions
and seek external evidence about them.

\citet{Goldberger1978} distinguished problems that are easily conflated.
The Honolulu model could not uniquely identify its parameters even from
population correlations; fixing a parameter could restore an estimable
version without establishing that the chosen value was correct. Separately,
misspecified models could produce wrong structural estimates despite good
fit. His simple twin example shows how excess environmental resemblance
among MZ twins can be attributed to heritability while the fitted equations
reproduce both twin correlations exactly \citep[p.~72]{Goldberger1978}.
The issue is not merely imprecision. It is what the observations establish
about the decomposition imposed on them.

\subsection{What Omitting a Model-Implied Restriction Does}
\citet[p.~19]{Goldberger1978} criticizes freeing spouse genetic resemblance while
retaining formulas derived under Fisher's mating scheme. Alternative
mating mechanisms may remove that restriction, but require a new derivation.
That is a warning about structural specification, not a theorem that an
unrestricted estimator must be inconsistent whenever a true equality is
left unimposed. If the unrestricted moments remain valid and uniquely
identify the true parameters, the true point still fits the population
moments exactly. Omission alone then does not change the probability limit.

Failing to impose a model-implied restriction can weaken or destroy
identification and can aggravate finite-sample bias. But persistent bias
under otherwise valid identification requires an additional argument,
such as incorrect moment equations or an invalid structural interpretation.
Goldberger's warning makes that investigation necessary; it does not
substitute for it. A statistically significant failure of a necessary
restriction challenges the joint maintained assumptions. It does not
identify which assumption failed or show that every estimate is biased.

Appendix~\ref{app:restriction-simulations} separates these cases using
population calculations and simulated pedigrees. Under the correct mating
model with identified unrestricted moments, both estimators recover the
population truth. With only parent--child and sibling moments, omitting
the phenotypic-assortment restriction instead leaves multiple triples
$(h^2,m,r)$ observationally equivalent. Under the wrong mating equations,
all three fitted parameters can converge to wrong values. Imposing the
restriction of that wrong model does not repair it.

\subsection{Confounding Is Not Limited to Genetic Nurture}
\label{sec:confounding-beyond-nurture}

Genetic nurture is one possible violation of the identifying exclusions,
not an exhaustive account of confounding. Shared environmental influences,
population structure, selection, or other unobserved determinants can
also generate covariance that the fitted kinship equations mistakenly
attribute to genetic transmission. To display the issue, write an outcome
on its original scale as $Y_i=G_i+U_i$, where $G_i$ is its direct genetic
contribution and $U_i$ collects the remaining influences. Then, for two
relatives,
\[
 \operatorname{Cov}(Y_i,Y_j)=\operatorname{Cov}(G_i,G_j)
 +\operatorname{Cov}(G_i,U_j)+\operatorname{Cov}(U_i,G_j)
 +\operatorname{Cov}(U_i,U_j).
\]
The maintained model specifies or excludes each of these terms; assortment
itself can generate some cross-person genetic--environmental covariance.
Any additional covariance omitted from the fitted equations can alter the
recovered genetic, assortment, and attenuation parameters. Rejecting a
particular parental-genotype pathway does not establish those exclusions.
An independent residual need not bias the genetic estimates, and different
violations need not bias every parameter or operate in the same direction:
the relevant question is whether the omitted influences change the moment
conditions used for identification.

This identification problem must also be distinguished from mediation.
A social response caused by the child's own genotype may be part of its
causal genetic effect, rather than confounding of that effect. Even a
correctly identified genetic effect therefore requires interpretation,
as Section~\ref{sec:interpretation} explains.


\subsection{A Perfect Fit Can Misclassify Cultural Transmission}
The distinction is especially clear in a deliberately constructed
counterexample. Let a socially transmitted characteristic follow the same
parental averaging process as an additive genetic characteristic, with
independent innovations and the same degree of assortment. Both are
Gaussian and independent of one another before matching. Their phenotypic
contributions then have the same kinship-correlation pattern. A model
that includes only genetic transmission absorbs their combined variance
into its genetic parameter.

In the example in Appendix~\ref{app:restriction-simulations}, the true
direct genetic variance share is $.48$, while a cultural component
contributes $.20$. The genetic-only model reports $h^2=.68$. It reproduces
all twelve fitted correlations exactly and satisfies its own restriction
$r=h^2m$. The construction is not offered as a realistic account of every
family process. It establishes that neither perfect fit nor satisfaction
of this restriction identifies the genetic interpretation. Independent
measurements or designs on which the mechanisms differ are required.
